% Part: computability
% Chapter: computability-theory
% Section: introduction

\documentclass[../../../include/open-logic-section]{subfiles}

\begin{document}

\olfileid{cmp}{thy}{int}
\olsection{پېژندنه}

د منطق هغه څانګه چې \emph{د محاسبه کېدنې تيوري} بلل کېږي، د تابعو او
سټونو د محاسبه کېدنې يا نسبي محاسبه کېدنې اړوندې مسئلې څېړي۔ دا د
کليني د اغېز نښه ده چې دې مضمون ته پخوا \emph{د بازګښت تيوري} ويل
کېدل، او نن دواړه نومونه په عام ډول کارېږي۔

تابع~$f\colon \Nat \pto \Nat$ هغه وخت \emph{جزوي محاسبه کېدونکې}
وبولئ چې د محاسبې په کوم مدل کښې محاسبه کېدے شي۔ که $f$ هرځاے تعريف
شوې وي، په ساده ډول به وايو چې $f$ \emph{محاسبه کېدونکې} ده۔ هغه
اړيکه~$R$ هم محاسبه کېدونکې بلل کېږي چې ځانګړونکې تابع~$\Char{R}$ ئې
محاسبه کېدونکې وي۔ که $f$ او~$g$ جزوي تابعې وي، $f(x) \fdefined$ به
په دې مانا ليکو چې $f$ پر~$x$ تعريف شوې ده، يعنې $x$ د~$f$ په تعريف
ساحه کښې دے؛ او $f(x) \fundefined$ به د دې مخالفه مانا، يعنې دا چې
$f$ پر~$x$ تعريف شوې نۀ ده، څرګندوي۔ $f(x) \simeq g(x)$ به په دې
مانا کاروو چې يا $f(x)$ او~$g(x)$ دواړه تعريف شوي نۀ دي، يا دواړه
تعريف شوي او سره برابر دي۔

د محاسبه کېدنې تيوري د محاسبې کوم ځانګړي مدل ته له مراجعې پرته هم
څېړل کېدے شي۔ د دې دپاره ښودل کېږي چې يوه نړيواله جزوي محاسبه
کېدونکې تابع~$\fn{Un}(k,x)$ شته۔ په دې سره جزوي محاسبه کېدونکې تابعې
په لړ کښې راوستل کېدے شي۔ موږ به $\cfind{k}$ د~$k$-مې يوځاييزې جزوي
محاسبه کېدونکې تابع دپاره وکاروو، چې د
$\cfind{k}(x) \simeq \fn{Un}(k,x)$ له لارې تعريفېږي۔ (کليني د دې
موخې دپاره $\{ k \}$ کارولی و، خو په وروستيو کښې دا نښه هومره نۀ ده
کارول شوې۔) لږ لا په عمومي ډول، د هرې ځاييزې کچې جزوي محاسبه کېدونکې
تابعې په يو شان ډول په لړ کښې راوستل کېدے شي، او $\cfind{k}[n]$ به د
$k$-مې $n$-ځاييزې جزوي محاسبه کېدونکې تابع دپاره کاروو۔
\textbf{د ژباړې د نسخې سمون (OLCMP-018):} سرچينې په همدې پراګراف کښې
لومړے نړيواله شمېرنه د جزوي محاسبه کېدونکو تابعو دپاره ټاکلې وه، خو
په وروستۍ فقره کښې ئې بې له کومې پېژندنې نوم جزوي بازګشتي تابعو ته
واړاوه۔ دلته نوم د پراګراف له ټاکلي، مدل-خپلواک ټولګي سره يو شان ساتل
شوے دے؛ نښه او ځاييزه کچه نۀ دي بدل شوي۔

که $f(\vec x,y)$ هرځاے تعريف شوې يا جزوي تابع وي، نو
$\umin{y}{f(\vec x,y)}$ د~$\vec x$ هغه تابع ده چې تر ټولو وړوکے~$y$
راګرځوي چې $f(\vec x,y)=0$ وي، په دې شرط چې $f(\vec x,0)$، \dots،
$f(\vec x,y-1)$ ټول تعريف شوي وي؛ که داسې~$y$ نۀ وي،
$\umin{y}{f(\vec x,y)}$ تعريف شوے نۀ دے۔ که $R(\vec x,y)$ يوه اړيکه
وي، $\umin{y}{R(\vec x,y)}$ تر ټولو وړوکے~$y$ تعريفېږي چې
$R(\vec x,y)$ رښتيا وي؛ په بله وينا، تر ټولو وړوکے~$y$ چې
$1 \tsub \Char{R}(\vec x,y)=0$ وي۔

د دې د ښودلو دپاره چې يوه تابع محاسبه کېدونکې ده، دوه لارې شته:
\begin{enumerate}
\item په کلک ډول: يو ټيورينګ ماشين يا جزوي بازګشتي تابع په څرګند ډول
  بيان کړئ، او وښيئ چې همدا موخه شوې تابع محاسبه کوي؛
\item په غير رسمي ډول: هغه الګوريتم بيان کړئ چې تابع محاسبه کوي، او د
  چرچ اصل ته مراجعه وکړئ۔
\end{enumerate}
د دواړو تر منځ تېره پوله نشته؛ د يو الګوريتم مفصل بيان بايد دومره
معلومات ورکړي چې په اصل کښې د سم ټيورينګ ماشين يا د جزوي بازګشتي
تعريفونو د لړۍ طرحه نسبتاً څرګنده وي۔ بشپړ کلک تعريفونه ښايي معلوماتي
نۀ وي، او هڅه به کوو چې د دغو دوو لارو تر منځ مناسب منځنی دريځ ومومو؛
په لنډ ډول، هڅه به کوو داسې شهودي خو کلک ثبوتونه ومومو چې ترې څرګنده
وي کره تعريفونه تر لاسه کېدے شي۔

\end{document}
